\documentclass[10pt, twocolumn, landscape, a4paper]{article}
% CEPRU | 3er Examen - Ciclo Ordinario 2017-I / Tema P | Grupo A..D
\usepackage[margin=1.5cm]{geometry}
\usepackage[utf8]{inputenc}
\usepackage[spanish,es-noshorthands]{babel}
\usepackage{amssymb, amsmath, amsbsy}
\usepackage{tasks}
\usepackage{enumerate}
\usepackage{graphicx} % \scalebox{1.4}{$$}
\usepackage[pdftex]{hyperref}
\hypersetup{pdftitle={},pdfauthor={Luis Carrillo Gutiérrez}}

\fontfamily{cmss}\selectfont
\renewcommand{\familydefault}{\sfdefault}
\pagestyle{empty} 
\setlength{\columnsep}{1.2in}


\begin{document}
\subsubsection*{Aritmética}
\begin{itemize}
\item{En las siguientes proposiciones, escribir (\textrm{V}) si es verdadera o (\textrm{F}) si es falsa. 
\begin{enumerate}[I.]
\item{En toda proporción geométrica, la suma de los extremos es igual a la suma de los medios.}
\item{Si : $\dfrac{a}{b} = \dfrac{c}{d} = k$, entonces $\dfrac{a + b}{b} = \dfrac{c + d}{d} = k$}
\item{Si : $\dfrac{a}{b} = \dfrac{c}{d}$, entonces $\dfrac{a + c}{b - c} = \dfrac{b + d}{b - d}$}
\item{En toda proporción geométrica continúa, el término medio es la semisuma de los términos extremos.}
\end{enumerate}
La secuencia correcta es :
\begin{tasks}(5)
\task{FFVF}\task{FVVF}\task{VVFF}\task{FFVV}\task{VFFV}
\end{tasks}
}
\item{Se desea repartir una herencia entre tres hermanos en forma directamente proporcional a sus edades que son : $17$, $13$ y $10$ años, pero si el reparto se hace un año después, el mayor quedaría perjudicado en $110$ soles. El monto de la herencia, en soles es :
\begin{tasks}(5)
\task{19 600}\task{18 200}\task{17 200}\task{14 400}\task{14 300}
\end{tasks}
}
\item{Si Denilson tarda $3$ horas en pintar una pared de forma cuadrada, entonces el número de horas que tardará en pintar otra pared de la misma forma cuyo lado es el triple del anterior, es :
\begin{tasks}(5)
\task{36}\task{27}\task{21}\task{9}\task{6}
\end{tasks}
}
\item{Para fijar el precio de un artículo se aumentó su costo en $300$ soles, pero al momento de venderse se hizo una rebaja del 10\% de este precio fijado, aún así se ganó el 10\% del precio de costo; entonces el precio fijado inicialmente, es : 
\begin{tasks}(3)
\task{1950 soles}
\task{1650 soles}
\task{1600 soles}
\task{1350 soles}
\task{1300 soles}
\end{tasks}
}
\item{Sofia observa que el interés producido por un capital en $5$ meses es $\dfrac1{10}$ del monto producido en este tiempo, ella desea saber qué parte del monto producido en $15$ meses será el interés producido en esos $15$ meses.
\begin{tasks}(5)
\task{$\dfrac1{5}$}
\task{$\dfrac{3}{10}$}
\task{$\dfrac1{4}$}
\task{$\dfrac{3}{4}$}
\task{$\dfrac1{3}$}
\end{tasks}
}
\item{En las siguientes proposiciones, escribir (\textrm{V}) si es verdadera o (\textrm{F}) si es falsa.
\begin{enumerate}[I.]
\item{El descuento racional ($\mathrm{D}_r$) es el interés simple que gana el valor actual comercial ($\mathrm{V}_{ac}$) de una letra de cambio.}
\item{Para los descuentos comercial ($\mathrm{D}_c$) y racional ($\mathrm{D}_r$) se verifica que : $\mathrm{D}_r - \mathrm{D}_c = \mathrm{V}_{ac} - \mathrm{V}_{ar}$, donde ($\mathrm{V}_{ac}$) es valor actual comercial y ($\mathrm{V}_{ar}$) es el valor actual racional.}
\item{La diferencia del descuento comercial ($\mathrm{D}_c$) y el descuento racional ($\mathrm{D}_r$) es el interés simple que gana el descuento racional ($\mathrm{D}_r$) de una letra de cambio.}
\item{Al descuento racional ($\mathrm{D}_r$) se le denomina también \emph{descuento matemático}.}
\end{enumerate}
La secuencia correcta es :
\begin{tasks}(5)
\task{VVVV}
\task{FVFF}
\task{FFVV}
\task{VFVF}
\task{FVVV}
\end{tasks}
}
\end{itemize}
\end{document}
